% Step 0.1: implementation details (environment and rasterizer check).
% Numbers from results/step0.1.json and results/step0.1_imagecas.json.
% Bib keys to add to the .bib: zha2024r2gaussian, biguri2016tigre, zeng2023imagecas.

\paragraph{Implementation details.}
We build on the official R$^2$-Gaussian code~\cite{zha2024r2gaussian}, whose
X-ray rasterizer and voxelizer we use without source modification, and on
TIGRE~\cite{biguri2016tigre} for cone-beam forward projection and FDK
initialization. All experiments run on a single NVIDIA RTX PRO 6000 Blackwell
GPU (compute capability 12.0) with PyTorch 2.11 and CUDA 13.0; the CUDA
extensions are compiled for this architecture with build flags only.
As a sanity check of this environment, R$^2$-Gaussian trained on its own
50-view synthetic chest case reaches a 3D PSNR of 35.15\,dB, against
35.09\,dB reported by the authors, in 13 minutes.
To confirm that the rasterizer handles thin vascular structures, we
reconstruct one coronary tree from ImageCAS~\cite{zeng2023imagecas}: the
largest connected component of the expert label of case 11 is resampled to a
$256^3$ volume of 0.37\,mm voxels, projected with the R$^2$-Gaussian cone-beam
protocol (50 views over $360^\circ$, Poisson and Gaussian noise), and
reconstructed with default settings. The result recovers the vessel with a
Dice score of 0.977 and a 3D PSNR of 45.17\,dB in 5.5 minutes.
